\section{DeSFAM: Dynamic eBPF-Driven Syscall Filtering and Anomaly Mitigation}
\label{sec:desfam}

This section presents the design and implementation of the proposed DeSFAM framework, detailing
its architecture, operational phases, and enforcement strategy.

\subsection{Overview}
\label{subsec:overview}

DeSFAM (Dynamic eBPF-Driven Syscall Filtering and Anomaly Mitigation) is a lightweight,
real-time framework designed to reduce the kernel attack surface and detect anomalous syscall
behavior in containerized environments. It enforces least-privilege syscall usage to mitigate
threats such as privilege escalation, container escape, and syscall injection.
Fig.~\ref{fig:architecture} illustrates the architectural workflow of the framework.

DeSFAM operates in three integrated phases:
\begin{itemize}
  \item \textbf{Hybrid Syscall Access Listing and Seccomp Profiling:} Combines static call graph
        extraction and dynamic eBPF tracing to construct a refined syscall access list, pruning
        unnecessary and CVE-associated syscalls.
  \item \textbf{Real-Time Monitoring and Anomaly Detection (SyscallAD):} Traces runtime syscall
        sequences and analyzes them using a VAE + Isolation Forest ensemble to detect
        sequence-level anomalies, including malicious use of permitted syscalls.
  \item \textbf{Adaptive Enforcement and Policy Management:} Enforces blocking, rate-limiting,
        or logging of high-risk syscalls using eBPF maps and LSM hooks, guided by risk scores
        from SyscallAD, without disrupting container execution.
\end{itemize}

By integrating sequence-aware profiling, anomaly detection, and real-time enforcement, DeSFAM
offers scalable, adaptive protection suited for modern cloud-native container deployments.

\begin{figure}[htbp]
  \centering
  \includegraphics[width=\linewidth]{figure1}
  \caption{DeSFAM architecture showing the interaction between user-space syscall profiling,
           anomaly detection, and kernel-space enforcement using eBPF and seccomp for adaptive,
           real-time syscall control.}
  \label{fig:architecture}
\end{figure}

\subsection{Design Considerations}
\label{subsec:design}

The design of DeSFAM is guided by the need to enforce syscall least privilege while addressing
the limitations of static-only filtering, high-overhead provenance tracking, and
retraining-dependent anomaly detectors. Traditional seccomp profiles are either too permissive
or brittle, failing to adapt to evolving runtime behavior. DeSFAM overcomes this by combining
static and dynamic analysis, enabling accurate and context-aware syscall profiling.

To ensure real-time enforcement without impacting container availability, DeSFAM leverages eBPF
probes and LSM hooks, which operate entirely in kernel space and support dynamic policy updates
with minimal latency. This design choice avoids the performance cost of continuous provenance
graph construction and eliminates the need for kernel modification or container restarts.

For anomaly detection, DeSFAM adopts a hybrid model (VAE + Isolation Forest) that captures
behavioral patterns in syscall sequences while avoiding frequent retraining, supporting scalable
and robust detection of zero-day attacks. Unlike hardware-bound solutions such as TEEs, DeSFAM
remains fully software-based and portable, ensuring compatibility across a wide range of
cloud-native deployments.

\subsection{Phase 1: Hybrid Syscall Access Listing and Seccomp Profiling}
\label{subsec:phase1}

Container security hinges on precise syscall control, as containers share the host kernel and
pose risks of unauthorized kernel access. Phase~1 of DeSFAM mitigates this by constructing a
container-specific syscall access list through a hybrid approach combining static analysis,
dynamic profiling, and CVE-based threat intelligence to generate a minimal and accurate seccomp
profile.

\begin{figure}[htbp]
  \centering
  \includegraphics[width=\linewidth]{figure2}
  \caption{DeSFAM Phase~1: Hybrid syscall access listing and seccomp profiling using static
           analysis, eBPF monitoring, and CVE-informed filtering to generate secure
           container-specific syscall policies.}
  \label{fig:phase1}
\end{figure}

\subsubsection{Static Analysis}
\label{subsubsec:static_analysis}

Static analysis begins by examining container images, including Dockerfiles, binaries, and
linked libraries, without executing them. Using tools like \texttt{objdump} and \texttt{readelf},
ELF binaries are disassembled to extract low-level syscall instructions and libc wrapper
functions.

To enhance coverage, a Static Call Graph (SCG) is constructed using NetworkX to trace
function-to-syscall mappings across the codebase. The resulting static syscall set,
$S_{\text{static}}$, captures all identifiable syscalls reachable from the code's control flow.

\begin{itemize}
  \item \textbf{Input:} Container image (Dockerfile, ELF binaries)
  \item \textbf{Tools:} \texttt{objdump}, \texttt{readelf}, \texttt{networkx}
  \item \textbf{Process:} Disassemble ELF binaries, extract syscall paths via SCG traversal
  \item \textbf{Output:} Static syscall list ($S_{\text{static}}$)
\end{itemize}

However, static analysis cannot detect syscalls invoked by runtime-resolved libraries or
JIT-compiled code. Thus, while $S_{\text{static}}$ offers a foundational syscall set, it must
be refined through dynamic analysis in subsequent phases.

\subsubsection{Dynamic Profiling: Capturing Runtime Behavior}
\label{subsubsec:dynamic_profiling}

To address the limitations of static analysis, DeSFAM incorporates dynamic profiling by
monitoring live syscall activity during container execution under standardized representative
workloads. This phase captures syscalls triggered by runtime conditions including dynamically
loaded libraries, environment-specific logic, and dependency interactions, while ensuring
comprehensive coverage of typical operational behaviors.

Containers execute under role-specific workloads designed using industry-standard benchmarks
that reflect real-world operational profiles:

\begin{itemize}
  \item \textbf{Web servers (Nginx):} \texttt{wrk} with 500--10,000 concurrent connections;
        ApacheBench simulating GET/POST requests at 1K--10K req/sec.
  \item \textbf{Databases (MySQL/Redis/MongoDB):} SysBench OLTP workloads (100--500 threads,
        1M+ transactions); \texttt{redis-benchmark} at 1M ops/sec.
  \item \textbf{Application runtimes (Node.js/Python):} Locust load tests simulating 50--500
        concurrent user flows; real-world dataset processing (GB-scale JSON/CSV parsing).
\end{itemize}

These workloads cover 90th-percentile production traffic patterns, container startup/shutdown
sequences, dependency interactions, and failure scenarios.

During execution, eBPF probes are attached to syscall entry points (\texttt{sys\_enter\_*}) to
trace syscall events. The collected metadata includes syscall IDs, process IDs, timestamps, and
namespace identifiers.

\begin{itemize}
  \item \textbf{Input:} Container instances under standardized workloads
  \item \textbf{Tools:} eBPF via \texttt{bcc}/\texttt{bpftrace} with kernel-space tracing
  \item \textbf{Process:} Deploy containers $\to$ attach eBPF probes $\to$ log metadata to ring
        buffers $\to$ aggregate syscalls across 3 workload cycles
  \item \textbf{Output:} Dynamic syscall list ($S_{\text{dynamic}}$)
\end{itemize}

\subsubsection{Threat Intelligence: CVE-Based Syscall Filtering}
\label{subsubsec:cve_filtering}

To enhance security awareness, DeSFAM integrates threat intelligence by incorporating CVE data
from the National Vulnerability Database (NVD). Using natural language processing (NLP), we
extract syscall names (e.g., \texttt{ptrace}, \texttt{clone}, \texttt{execve}) from CVE
descriptions to build a list of high-risk syscalls, denoted as $S_{\text{blocked}}$.

To construct $S_{\text{blocked}}$, we employ an NLP-based parsing engine that analyzes CVE
reports retrieved from the NVD via its REST API. A rule-based keyword extraction technique is
applied to identify syscall names mentioned in CVE descriptions. We utilize a curated dictionary
of Linux syscall names derived from kernel source headers (e.g., \texttt{unistd64.h}) to ensure
accurate mapping. Furthermore, to avoid overblocking, contextual metadata such as kernel version
applicability, affected subsystems, exploit conditions, and argument constraints are extracted
from the CVE text using pattern matching and regular expression parsing.

\begin{itemize}
  \item \textbf{Input:} CVE entries with syscall mentions
  \item \textbf{Mapping:} Rule-based syscall extraction + context filtering
  \item \textbf{Output:} Blocked syscall list ($S_{\text{blocked}}$)
\end{itemize}

This set is excluded from the union of static and dynamic syscall sets, ensuring vulnerable
syscalls are omitted even if observed during profiling.

\subsubsection{Final Access List Synthesis and Policy Generation}
\label{subsubsec:synthesis}

The final access list $S_{\text{final}}$ is computed as:
\begin{equation}
  S_{\text{final}} = \bigl(S_{\text{static}} \cup S_{\text{dynamic}}\bigr)
                      \cap S_{\text{template}}
                      \setminus S_{\text{blocked}}
                      \setminus S_{\text{template}}
  \label{eq:sfinal}
\end{equation}

Here, $S_{\text{template}}$ defines role-specific expected syscall behavior for the container
(e.g., web server, database, GPU workload), serving as a functional boundary. The template
syscalls are derived by combining several sources: official Docker default seccomp profiles,
Linux kernel syscall documentation, community-maintained security rules from the Falco project,
and empirical syscall traces collected during initial benign container executions.

This multi-source approach ensures that rarely invoked but necessary syscalls during container
initialization or uncommon execution paths are incorporated into $S_{\text{template}}$ while
maintaining minimal syscall exposure.

The CVE-based syscall set $S_{\text{blocked}}$ is generated by extracting exploit conditions,
including kernel versions, device-specific vulnerabilities, or specific syscall argument
patterns. This prevents overblocking syscalls that are only vulnerable under specific
configurations not applicable to the container's runtime environment.

The resulting access list $S_{\text{final}}$ is compiled into a seccomp policy
(\texttt{seccomp\_profile.json}) enforced by the Linux kernel at runtime. Fig.~\ref{fig:phase1}
and Algorithm~\ref{alg:phase1} summarize this process.

\begin{algorithm}[htbp]
\caption{Enhanced Hybrid Syscall Profiling}
\label{alg:phase1}
\begin{algorithmic}[1]
  \Require Container image, running instance, CVE database
  \Ensure Secure, functional syscall access list
  \State Extract syscalls via static analysis $\to S_{\text{static}}$
  \State Profile runtime under workload (dynamic profiling) $\to S_{\text{dynamic}}$
  \State Parse CVEs with contextual analysis $\to S_{\text{blocked}}$
  \State Generate $S_{\text{final}} = (S_{\text{static}} \cup S_{\text{dynamic}})
         \cap S_{\text{template}} \setminus S_{\text{blocked}} \setminus S_{\text{template}}$
  \State \Return $S_{\text{final}}$
\end{algorithmic}
\end{algorithm}

\subsubsection{Container-Specific Policy Upload}
\label{subsubsec:policy_upload}

DeSFAM dynamically loads container-specific syscall policies using \texttt{policy\_loader}.
Each policy reflects the container's behavioral profile from prior analysis and is written into
the pinned eBPF map at \texttt{/sys/fs/bpf/blocked\_syscalls}. Namespace identifiers serve as
keys for container isolation.

This design enforces least-privilege syscall access per container at runtime, without restarting
instances. Fig.~\ref{fig:policy_upload} shows the upload pipeline and the population of the
Syscall map.

\begin{figure}[htbp]
  \centering
  \includegraphics[width=\linewidth]{figure3}
  \caption{DeSFAM Phase~3: (a) Dynamic upload of container-specific syscall policies using
           \texttt{policy\_loader}; (b) Verification of blocked syscalls in the pinned eBPF
           map using \texttt{bpftool}.}
  \label{fig:policy_upload}
\end{figure}

\subsection{Phase 2: Real-Time Monitoring and Syscall Anomaly Detection (SyscallAD)}
\label{subsec:phase2}

Phase~2 of DeSFAM introduces SyscallAD, a real-time anomaly detection module that combines deep
learning and ensemble methods to identify both known and zero-day threats in containerized
environments. It monitors syscall sequences using eBPF and detects behavioral deviations linked
to privilege escalation, container escape, or syscall injection.

SyscallAD employs a hybrid ensemble comprising a Variational Autoencoder (VAE), which learns
latent representations of benign syscall patterns, and an Isolation Forest (iForest), which
identifies statistical outliers. This combination enables robust, low-latency detection with
minimal false positives and overhead.

\begin{figure}[htbp]
  \centering
  \includegraphics[width=\linewidth]{figure4}
  \caption{DeSFAM Phase~2: Real-time syscall monitoring and anomaly detection using a hybrid
           VAE + iForest model. The offline module is trained on DongTing sequences, while the
           online module processes live syscalls via eBPF for anomaly evaluation and policy
           escalation.}
  \label{fig:phase2}
\end{figure}

\subsubsection{The DongTing Dataset: Behavioral Baselines for Detection}
\label{subsubsec:dongting}

SyscallAD is trained on the DongTing dataset~\cite{ref26_dongtingdataset}, which comprises
18,966 syscall sequences and over 2.6~million individual syscalls collected from containerized
applications. The dataset is split into:

\begin{itemize}
  \item \textbf{6,850 benign sequences:} Captured from normal workloads such as web servers and
        databases executing file I/O (\texttt{open}, \texttt{read}) and memory operations
        (\texttt{mmap}, \texttt{futex}).
  \item \textbf{12,116 malicious sequences:} Generated through simulated attacks, including
        privilege escalation (\texttt{setuid}), filesystem tampering (\texttt{unlink}), and
        network intrusion (\texttt{socket}).
\end{itemize}

Each sequence includes syscall IDs, timestamps, and contextual metadata (e.g., process IDs,
namespaces). Given the inherent class imbalance, sampling techniques were applied to prevent
model bias.

\subsubsection{PrefixSpan Pattern Mining: Extracting Trusted Syscall Templates}
\label{subsubsec:prefixspan}

To reduce false positives and improve detection efficiency, DeSFAM integrates
PrefixSpan~\cite{prefixspan_pei}---a sequential pattern mining algorithm---on the benign subset
of the DongTing dataset. This identifies frequent syscall subsequences that represent legitimate
container behavior.

The resulting Benign Pattern Database is used to:
\begin{itemize}
  \item Match incoming syscall windows against known-good patterns for early benign
        classification.
  \item Bypass deep anomaly scoring for matched sequences, reducing computational overhead.
  \item Provide semantic traceability for security decisions.
\end{itemize}

\subsubsection{Feature Engineering: Structuring Syscall Sequences for Detection}
\label{subsubsec:features}

SyscallAD applies multi-dimensional feature engineering to convert raw syscall streams into
actionable model inputs:

\begin{itemize}
  \item \textbf{Sliding Windows:} Syscalls are segmented into overlapping windows (length~15,
        stride~3), capturing both short-term anomalies and extended behavior patterns.
  \item \textbf{Categorical Frequencies:} Syscalls are mapped to functional categories (e.g.,
        file I/O, memory, network), and their frequency per window forms a behavioral signature.
  \item \textbf{Temporal Features:} Inter-syscall time intervals (temporal deltas) are
        summarized (mean, std, max) to detect abnormal timing behaviors.
  \item \textbf{Access List Pattern Matching:} Each window is cross-checked against benign
        patterns from PrefixSpan; matches increase confidence in benign classification, reducing
        false positives.
  \item \textbf{Data Augmentation:} Malicious sequences are synthetically modified (e.g.,
        insertion or permutation of high-risk syscalls like \texttt{clone}, \texttt{setuid}) to
        simulate zero-day scenarios, improving model robustness.
\end{itemize}

\subsubsection{SyscallAD: Hybrid Deep Learning and Anomaly Detection Engine}
\label{subsubsec:syscallad}

SyscallAD is a hybrid anomaly detection module designed to identify behavioral deviations in
syscall sequences with high precision and adaptability.

\paragraph{Optimized Variational Autoencoder (VAE)}
The VAE employs an encoder-decoder architecture ($32 \to 8 \to 32$ neurons) with SELU
activations, dropout ($p = 0.2$), and L2 regularization. It learns to reconstruct benign
syscall windows by minimizing Mean Squared Reconstruction Error (MSRE). Anomalies are flagged
when MSRE exceeds the 99.5th percentile of benign training errors.

\paragraph{Isolation Forest (iForest)}
iForest isolates anomalies through recursive partitioning. Rare syscall sequences yield shallower
isolation trees and higher anomaly scores. A contamination parameter of 2\% controls sensitivity,
allowing detection of statistically rare behaviors missed by VAE.

\paragraph{Ensemble Fusion Strategy}
SyscallAD combines the strengths of VAE and Isolation Forest through a weighted score fusion
mechanism. Each syscall window $W_i$ is evaluated by both models:
\begin{itemize}
  \item $A_{\text{VAE}}(W_i)$: Anomaly score based on Mean Squared Reconstruction Error.
  \item $A_{\text{iForest}}(W_i)$: Score derived from isolation depth in the Isolation Forest.
\end{itemize}

The final score is computed as:
\begin{equation}
  A(W_i) = \alpha \cdot A_{\text{VAE}}(W_i) + (1 - \alpha) \cdot A_{\text{iForest}}(W_i)
  \label{eq:ensemble}
\end{equation}
where $\alpha = 0.7$ prioritizes the VAE's contribution. This value was validated via
sensitivity analysis over $\alpha \in \{0.0, 0.25, 0.5, 0.7, 0.85, 1.0\}$ on the DongTing
dataset, as shown in Fig.~\ref{fig:sensitivity}.

\begin{figure*}[tp]
  \centering
  \includegraphics[width=\linewidth]{figure6}
  \caption{Sensitivity analysis of fusion weight $\alpha$ on F1-Score and AUC.
           Peak performance is observed at $\alpha = 0.7$.}
  \label{fig:sensitivity}
\end{figure*}

Fig.~\ref{fig:syscallad_output} presents the anomaly logs from \texttt{test-mysql} and
\texttt{test-redis} containers. Benign syscalls such as \texttt{futex}, \texttt{epoll\_wait},
and \texttt{wait4} are correctly classified, while anomalies like \texttt{clone} and
\texttt{ioctl}, indicative of escape or privilege escalation attempts, are effectively flagged.

\begin{figure*}[tp]
  \centering
  \includegraphics[width=\linewidth]{figure5}
  \caption{SyscallAD output from DeSFAM Phase~2: normal behavior in \texttt{test-mysql}
           contrasted with anomalous syscall detections in \texttt{test-redis}
           (e.g., \texttt{clone}, \texttt{ioctl}).}
  \label{fig:syscallad_output}
\end{figure*}

\paragraph{Dynamic Thresholding: Adaptive Anomaly Boundary Adjustment}
To accommodate evolving syscall behavior, SyscallAD employs dynamic thresholding based on an
Exponential Moving Average (EMA), rather than a fixed cut-off. The initial threshold $T_0$ is
set to the 99.5th percentile of VAE reconstruction errors on benign validation data. During
runtime, the threshold is updated after each syscall window $W_i$ as:
\begin{equation}
  T_{t+1} = \max\bigl(T_{\min},\; \beta T_t + (1 - \beta) A(W_i)\bigr)
  \label{eq:threshold}
\end{equation}
where $\beta = 0.9$ controls update sensitivity and $T_{\min} = 0.5$ enforces a lower bound.

Final detection is defined as:
\begin{equation}
  \mathrm{Flag}(W_i) =
  \begin{cases}
    1, & \text{if } A(W_i) > T \\
    0, & \text{otherwise}
  \end{cases}
  \label{eq:flag}
\end{equation}

\subsubsection{Real-Time Monitoring with eBPF: Low-Overhead Container-Aware Syscall Capture}
\label{subsubsec:ebpf_monitoring}

DeSFAM utilizes eBPF tracepoints to enable fine-grained, low-overhead syscall surveillance
directly at the Linux kernel level. The BPF program hooks into
\texttt{(raw\_syscalls:sys\_enter)} and \texttt{(raw\_syscalls:sys\_exit)}, enabling the
interception of syscall invocations at their entry and exit points.

To isolate syscall activity from containerized workloads, DeSFAM dynamically maps container
PIDs by querying the Docker runtime and injecting them into a \texttt{pid\_map} BPF hash table.
Only syscalls originating from these PIDs are recorded, ensuring container-specific tracing
without unnecessary system noise.

Captured syscall events are enriched with contextual metadata including: syscall ID, timestamp
(nanosecond resolution), process and parent PIDs, command name, container name, and syscall
return value. These are sent from kernel to user space via a \texttt{ring\_buffer} and logged to
per-container files in JSON format. Simultaneously, syscall IDs are appended to a rolling deque
(sliding window of size~15) in memory for each container. Fig.~\ref{fig:ebpf_log} shows a
representative real-time syscall log for a MySQL container.

\begin{figure*}[tp]
  \centering
  \includegraphics[width=\linewidth]{figure7}
  \caption{Real-time eBPF syscall log for \texttt{kubepro-mysql-1}, showing timestamped syscall
           events with PID, container name, user, and return value.}
  \label{fig:ebpf_log}
\end{figure*}

\subsubsection{DevOps Integration: Continuous Model Refinement}
\label{subsubsec:devops}

To mitigate false positives and adapt to evolving workloads, DeSFAM integrates a DevOps-driven
feedback loop within Phase~2. Anomalous syscall sequences, along with metadata (e.g., container
name, PID, timestamp), are logged and visualized for analyst review. Security teams can annotate
detections as true or false positives, with validated false positives appended to the benign
training set.

SyscallAD periodically retrains its VAE + iForest ensemble on this augmented dataset, enabling
it to learn new benign variations introduced by environment-specific workloads or edge cases.

The full detection workflow is formalized in Algorithm~\ref{alg:phase2}.

\begin{algorithm}[htbp]
\caption{Real-Time Syscall Anomaly Detection (SyscallAD)}
\label{alg:phase2}
\begin{algorithmic}[1]
  \Require eBPF hook, trained VAE and iForest models,
           validation reconstruction errors $E_{\text{val}}$
  \Ensure Anomaly flag for each syscall window
  \State Initialize sliding windows $W_i \leftarrow \emptyset$ for $i \in \{5, 10, 15\}$
  \State Load trained VAE and iForest models
  \State Compute $T_0 \leftarrow \text{percentile}(E_{\text{val}}, 99.5)$
  \State Set threshold $T \leftarrow T_0$
  \While{syscall stream is active}
    \State Capture syscall $s_t$ at time $t$
    \State Append $s_t$ to each $W_i$
    \For{each full window $W_i$}
      \State $f_i \leftarrow \text{extract\_features}(W_i)$
      \State $A_{\text{VAE}} \leftarrow \text{VAE\_score}(f_i)$
      \State $A_{\text{iForest}} \leftarrow \text{iForest\_score}(f_i)$
      \State $A(W_i) \leftarrow 0.7 \cdot A_{\text{VAE}} + 0.3 \cdot A_{\text{iForest}}$
      \If{$A(W_i) > T$}
        \State Flag $W_i$ as anomalous; escalate to Phase~3
      \Else
        \State $T \leftarrow \max(T_{\min},\; \beta \cdot T + (1 - \beta) \cdot A(W_i))$
      \EndIf
      \State Slide $W_i$ forward
    \EndFor
  \EndWhile
\end{algorithmic}
\end{algorithm}

\subsection{Phase 3: Adaptive Syscall Enforcement and Policy Management}
\label{subsec:phase3}

Adaptive Syscall Enforcement is the third phase of the DeSFAM framework. It focuses on
dynamically enforcing syscall policies in containerized environments in real time. This phase
integrates risk scoring, contextual analysis, and external threat intelligence to deliver
adaptive and precise security responses.

\subsubsection{Risk Scoring Mechanism}
\label{subsubsec:risk_scoring}

DeSFAM quantifies the threat level of each anomalous syscall sequence $W_i$ using a composite
risk score derived from three components.

\paragraph{MITRE ATT\&CK Mapping}
Each syscall in $W_i$ is mapped to associated adversarial techniques (e.g., T1055 for
\texttt{ptrace}) based on a predefined lookup (Table~\ref{tab:mitre_mapping}). This contextual
mapping enables reasoning beyond isolated calls to recognize high-level attack patterns.

\begin{table*}[tp]
  \centering
  \caption{Example Syscall-to-MITRE ATT\&CK Technique Mapping.}
  \label{tab:mitre_mapping}
  \includegraphics[width=\linewidth]{table2}
\end{table*}

\paragraph{Cumulative Severity}
The cumulative severity score aggregates the individual severity scores of syscalls in the
sequence:
\begin{equation}
  \text{Cumulative Severity} = \sum_{t \in W_i} \text{Severity}(t)
  \label{eq:severity}
\end{equation}

\paragraph{Exploit Pattern Matching}
Each sequence $W_i$ is evaluated against known CVE exploit signatures. If a match is found, an
additional risk penalty is applied:
\begin{equation}
  I(\text{Exploit Match}) =
  \begin{cases}
    1, & \text{if matched} \\
    0, & \text{otherwise}
  \end{cases}
  \label{eq:exploit_match}
\end{equation}

\paragraph{Final Risk Score}
The final risk score integrates both severity and exploit intelligence:
\begin{equation}
  \text{Risk Score} = \sum_{t \in W_i} \text{Severity}(t) + 5 \cdot I(\text{Exploit Match})
  \label{eq:risk_score}
\end{equation}

\paragraph{Policy Decision Logic}
Based on the computed risk score, DeSFAM applies a three-tier enforcement strategy:
\begin{equation}
  \text{Action} =
  \begin{cases}
    \text{Block},       & \text{if Risk Score} \geq 8 \\
    \text{Rate-limit},  & \text{if } 5 \leq \text{Risk Score} < 8 \\
    \text{Log},         & \text{if Risk Score} < 5
  \end{cases}
  \label{eq:action}
\end{equation}

These thresholds are empirically tuned and remain configurable for deployment-specific risk
postures. Table~\ref{tab:risk_examples} presents sample sequences and corresponding actions.

\begin{table}[htbp]
  \centering
  \caption{Example Risk Scores for Benign and Malicious Sequences.}
  \label{tab:risk_examples}
  \includegraphics[width=\linewidth]{table3}
\end{table}

\subsubsection{eBPF-Based Enforcement}
\label{subsubsec:ebpf_enforcement}

DeSFAM enforces syscall policies through eBPF programs attached to LSM hooks. Specifically,
these programs are loaded and attached using \texttt{bpf\_program\_attach\_lsm()} to intercept
system calls at the LSM's \texttt{security\_syscall} hook.

\paragraph{LSM Hook Attachment}
eBPF programs are dynamically attached to the \texttt{security\_syscall} LSM hook:

\begin{lstlisting}[float=htbp, language=C, caption={eBPF LSM hook attachment.}]
// eBPF Hook Attachment
struct bpf_program *prog =
    bpf_object__find_program_by_name(obj, "enforce_policy");
struct bpf_link *link = bpf_program_attach_lsm(prog);
\end{lstlisting}

The hook function signature:
\begin{lstlisting}[float=htbp, language=C, caption={LSM hook function signature.}]
SEC("lsm/syscall")
int BPF_PROG(enforce_policy, int syscall_num,
             struct pt_regs *regs) {
    // Enforcement logic here
}
\end{lstlisting}

\paragraph{Policy Map Structure}
Syscall actions are stored in a hash-based eBPF map indexed by container identifiers:

\begin{lstlisting}[float=htbp, language=C, caption={Container-keyed eBPF policy map.}]
struct container_key {
    u32 mntns;  // Mount namespace ID
    u32 pidns;  // PID namespace ID
};

struct bpf_map_def SEC("maps") policy_map = {
    .type       = BPF_MAP_TYPE_HASH,
    .key_size   = sizeof(struct container_key),
    .value_size = sizeof(u64),
    .max_entries = 144,
};
\end{lstlisting}

The map is pinned at \texttt{/sys/fs/bpf/blocked\_syscalls} for runtime updates via
\texttt{policy\_loader}.

\paragraph{Syscall Blocking Mechanism}
The actual blocking of syscalls occurs within the eBPF hook attached to the
\texttt{security\_syscall} LSM. The critical enforcement primitive uses kernel function override:
\begin{lstlisting}[float=htbp, language=C, caption={Syscall blocking via kernel function override.}]
// Immediate syscall denial
bpf_override_return(regs, -EPERM);
\end{lstlisting}

The enforcement workflow:
\begin{enumerate}
  \item Retrieve container context: \texttt{mntns = bpf\_get\_ns\_id(NS\_MNT);
        pidns = bpf\_get\_ns\_id(NS\_PID);}
  \item Query policy map: \texttt{u64 *action = bpf\_map\_lookup\_elem(\&policy\_map, \&key);}
  \item Enforce action:
        \textbf{BLOCK}: Deny via \texttt{bpf\_override\_return(regs, -EPERM)};
        \textbf{RATE\_LIMIT}: Allow $N$ invocations/minute via atomic counters;
        \textbf{LOG}: Record metadata via ring buffer.
\end{enumerate}

\paragraph{Self-Healing Access List via Seccomp Regeneration}
To ensure long-term policy evolution, DeSFAM incorporates a self-healing mechanism that
automatically refines static syscall access lists based on repeated runtime blocks.

Blocked syscalls are tracked in-kernel using a Least Recently Used (LRU) cache map called
\texttt{probation\_map}:

\begin{lstlisting}[float=htbp, language=C, caption={Probation map for tracking repeated blocks.}]
struct probation_key {
    struct container_key ckey;
    u32 syscall;  // Syscall number
};

SEC("maps") struct bpf_map_def probation_map = {
    .type       = BPF_MAP_TYPE_LRU_HASH,
    .key_size   = sizeof(struct probation_key),
    .value_size = sizeof(u32),  // Block counter
    .max_entries = 512,
};
\end{lstlisting}

If a syscall is blocked three or more times within a configurable window and is not previously
annotated as benign, the kernel triggers a user-space daemon via
\texttt{bpf\_send\_signal(SIGUSR1)}. This daemon resolves the syscall name, locates and updates
the corresponding container's seccomp profile, and removes the flagged syscall from the access
list. This adaptive defense loop ensures that newly deployed containers inherit hardened,
exploit-aware syscall policies without requiring restarts or manual review.

\paragraph{Formal Policy Model}
Let $S$ be the set of valid syscalls, and $M$, $N$ represent mount and PID namespaces. The
enforcement policy is defined as:
\begin{equation}
  P : S \times M \times N \to \{0, 1\}
  \label{eq:policy_func}
\end{equation}
\begin{equation}
  \text{Decision}(s) =
  \begin{cases}
    \text{Allow}, & \text{if } P(s, \mathit{mntns}, \mathit{pidns}) = 1 \\
    \text{Block},  & \text{otherwise}
  \end{cases}
  \label{eq:decision}
\end{equation}

This container-aware logic ensures syscall isolation and enforces least-privilege access per
namespace. The combined use of eBPF tracepoints and LSM hooks enables fast, layered enforcement
with minimal overhead.

The enforcement procedure in Phase~3 is formalized through Algorithm~\ref{alg:phase3}, and the
complete enforcement pipeline is illustrated in Fig.~\ref{fig:phase3}.

\begin{figure}[htbp]
  \centering
  \includegraphics[width=\linewidth]{figure8}
  \caption{DeSFAM Phase~3 -- Adaptive enforcement pipeline. Anomalous syscall sequences from
           Phase~2 are intercepted and acted upon via eBPF and LSM hooks.}
  \label{fig:phase3}
\end{figure}

\begin{algorithm}[htbp]
\caption{Adaptive Syscall Enforcement With Auto-Hardening}
\label{alg:phase3}
\begin{algorithmic}[1]
  \Require Anomalous sequence $W_i$, MITRE mapping, exploit database,
           container context (\textit{mntns}, \textit{pidns}), eBPF map \texttt{policy\_map}
  \Ensure Enforcement action: ``Block'', ``Rate-limit'', or ``Log''
  \State Initialize $\mathit{Risk\_Score} \leftarrow 0$
  \For{each syscall $t$ in $W_i$}
    \State $\mathit{Risk\_Score} \leftarrow \mathit{Risk\_Score} + \text{Severity}(t)$
  \EndFor
  \If{$W_i$ matches an exploit in the database}
    \State $\mathit{Risk\_Score} \leftarrow \mathit{Risk\_Score} + 5$
  \EndIf
  \If{$\mathit{Risk\_Score} \geq 8$}
    \State $\mathit{action} \leftarrow$ ``Block''
  \ElsIf{$5 \leq \mathit{Risk\_Score} < 8$}
    \State $\mathit{action} \leftarrow$ ``Rate-limit''
  \Else
    \State $\mathit{action} \leftarrow$ ``Log''
  \EndIf
  \For{each syscall $t$ in $W_i$}
    \State Update $\texttt{policy\_map}[(t, \mathit{mntns}, \mathit{pidns})] \leftarrow \mathit{action}$
    \State Enforce $\mathit{action}$ on $t$ via eBPF
    \State Log event $(t, \mathit{action}, \mathit{mntns}, \mathit{pidns})$
  \EndFor
  \If{$\mathit{action} = $ ``Block''}
    \For{each syscall $t$ in $W_i$}
      \State $\mathit{probation\_count}[t] \leftarrow \mathit{probation\_count}[t] + 1$
      \If{$\mathit{probation\_count}[t] \geq 3$ \textbf{and} $t$ not annotated as benign}
        \State Trigger seccomp access list update for syscall $t$
      \EndIf
    \EndFor
  \EndIf
  \State \Return $\mathit{action}$
\end{algorithmic}
\end{algorithm}

\subsubsection{Use Case: Privilege Escalation Attempt and Real-World CVE Defense}
\label{subsubsec:usecase}

Consider a container running under DeSFAM's restrictive policy designed to block privilege
escalation attempts. A typical attack sequence---\texttt{clone()}, \texttt{setuid()},
\texttt{execve()}---is intercepted at the kernel level.

Assume the process operates in mount namespace $\mathit{mntns} = 1001$ and PID namespace
$\mathit{pidns} = 2001$. Upon invocation of \texttt{setuid()}, DeSFAM consults its
per-container policy: \texttt{policy\_map[1001, 2001] = [..., 1, 0, 1, ...]}, where each entry
denotes the action for a syscall ID ($0 = \text{BLOCK}$, $1 = \text{ALLOW}$). Since
\texttt{setuid()} is mapped to 0, the syscall is denied. The eBPF LSM hook,
\texttt{realtime\_blocking}, enforces this decision in the kernel, returning \texttt{-EPERM}
and logging the event for auditing. This ensures immediate, in-kernel mitigation of exploit
attempts without requiring container restarts or user-space intervention.
